\documentclass{Tesis}
\usepackage[utf8]{inputenc}
\usepackage[T1]{fontenc}
\usepackage{geometry}
\usepackage{graphicx}
\usepackage{float}
\usepackage{hyperref}
\usepackage[english,indonesian]{babel}

\title{Optimasi Pemrosesan Template LaTeX Menggunakan Pendekatan Berbasis AI}
\author{Academic Support Agent}
\date{\today}

\begin{document}
\maketitle

\begin{abstract}

\end{abstract}

\chapter{PENDAHULUAN}


\chapter{TINJAUAN PUSTAKA}


% Original content (commented out):
% \documentclass[12pt]{Tesis}
\usepackage{lipsum}
% Para links sin decoración
%\hypersetup{colorlinks=false,allbordercolors=white}

\title{La cuadratura del círculo}
\author{Carl Friedrich Gauss}
\grado{Maestro en Ciencias en Ingeniería Mecatrónica}
\date{Junio de 2018}
\director{Albert Einstein}
\codirector{Max Planck} % Si existe un codirector
%\tutor{Erwin Schrödinger} % Si existe un tutor
\autorizacion{autorizacion} % autorización de impresión digitalizada 
\DeclareGraphicsExtensions{.pdf,.png,.jpg}
\graphicspath{{./figuras/}}

\begin{document}

\maketitle

\frontmatter

\insertarautorizacion

\begin{dedicatoria}
\begin{center}
Escriba aquí su dedicatoria.
\end{center}
\end{dedicatoria}

\begin{abstract}
\lipsum[1-2]
\end{abstract}

\selectlanguage{english}
\begin{abstract}
\lipsum[1-2]
\end{abstract}
\selectlanguage{spanish}

\begin{agradecimientos}
\begin{center}
\textbf{AGRADECIMIENTOS}
\end{center}
Puede agradecer a quienes le asesoraron académicamente, a las instituciones que financiaron sus estudios, al sancho que cuidó a su novia mientras usted estudiaba, etc.

\lipsum[2]
\end{agradecimientos}

\tableofcontents
{\noskip\listoffigures}
{\noskip\listoftables}

\mainmatter

\chapter{Introducción}
\epigrafe{El todo es más que la suma de sus partes.}{\textsc{Aristóteles}}

\section{Primera sección}
\lipsum[1]
\begin{figure}[htbp]
\centering
\includegraphics[width=0.35\textwidth]{dron}
\caption{Vehículo aéreo no tripulado.}
\label{fig:dron}
\end{figure}
\lipsum[2]

\section{Segunda sección}
\lipsum[3-4]
En el \href{http://www.microsoft.com}{sitio web de Microsoft} se pueden consultar las características del vehículo mostrado en la Figura~\ref{fig:dron}. 

\section{Tercera sección}
\lipsum[5-6]

\section{Cuarta sección}
\lipsum[7-8]

\chapter{Fundamento teórico}

\section{Primera sección}

\lipsum[1]

\begin{equation}                                                     
\mathrm{e}^{x} = \sum_{k=0}^{\infty} \frac{x^k}{k!} = 
1 + x + \frac{x^2}{2} + \frac{x^3}{6} + \cdots
\label{eq:taylor}
\end{equation}

La ecuación \eqref{eq:taylor} muestra la expansión en serie de Taylor alrededor de $x = 0$ para la función exponencial natural.

\section{Segunda sección}
\lipsum[3-4]

\section{Tercera sección}
\lipsum[5-6]

\section{Cuarta sección}
\lipsum[7-8]

\chapter{Metodología}

\section{Primera sección}
\lipsum[1-2]

\section{Segunda sección}
\lipsum[3]
\begin{figure}[htbp]
\centering
\includegraphics{particion}
\caption{Partición difusa de la temperatura en el reactor.}
\label{fig:particion}
\end{figure}
\lipsum[4]

\section{Tercera sección}
\lipsum[5-6]

\section{Cuarta sección}
\lipsum[7-8]

\chapter{Resultados}

\section{Primera sección}
\lipsum[1]
\begin{table}[htbp]
\centering
\caption{Variación de la presión con la temperatura.}
\label{tab:pversust}
\begin{tabular}{ccc}
\toprule
\textbf{Temperatura (\si{\kelvin})} & \textbf{Presión (\si{\kilo\pascal})} \\
\midrule
1.234 & 2.345 \\
3.456 & 4.567 \\ 
5.678 & 6.789 \\
7.890 & 8.901 \\
\bottomrule
\end{tabular}
\end{table}
\lipsum[2]

\section{Segunda sección}
\lipsum[3-4]

\section{Tercera sección}
\lipsum[5-6]

\section{Cuarta sección}
\lipsum[7-8]

\chapter{Conclusiones}
\lipsum[1-5]

\appendix

\chapter{Demostraciones}

\section{Teorema X}
\lipsum[1-2]

\section{Teorema Y}
\lipsum[1-2]


\chapter{Diagramas}

\section{Diagrama X}
\lipsum[1-2]

\section{Diagrama Y}
\lipsum[1-2]

\backmatter

\bibliographystyle{apacite}
\bibliography{referencias}
\nocite{*}

\end{document}

\end{document}